\documentclass[main]{subfiles}
\usepackage[sols]{workbook}
\numbering{ws2-13}

\begin{document}

\ifSubfilesClassLoaded{%
  \chaptermark{\chaptertwo}
}{}

\section{Series Wrap-Up} \label{ws2-13}
\begin{enumerate}

\item Series toolbox final review.
\begin{enumerate}
    \item 
    \begin{problem}[2.5in]
    Make a list of all the tests and criteria for determining convergence or divergence of series. Specify the conditions for being able to apply each of them.    
    \end{problem}

    \begin{solution}

    Here is one possible way of organizing the tests:

\begin{enumerate}
    \item \textbf{Tests based on type of series:}
        \begin{itemize}
            \item $p$-series
            \begin{itemize}
                \item \textit{Condition:} Only applies to series of the form $\displaystyle\sum_{k=1}^\infty \frac{1}{k^p}$

                \item \textit{Conclusion:} The series converges when $p > 1$ and diverges when $p \leq 1$
            \end{itemize}
            
            \item Geometric series
            \begin{itemize}
                \item \textit{Condition:} Only applies to series of the form $\displaystyle\sum_{k=0}^\infty ar^k$

                \item \textit{Conclusion:} The series converges to $a/(1-r)$ when $|r| <1$ and diverges when $|r| \geq 1$
            \end{itemize}

            \item Alternating series test
            \begin{itemize}
                \item \textit{Condition:} Only applies to series of the form $\displaystyle\sum_{k=1}^\infty (-1)^k|a_k|$

                \item \textit{Conclusion:} If $|a_{k+1}|<|a_k|$ and $\displaystyle\lim_{k \to \infty} |a_k| = 0$ then the series converges
            \end{itemize}
        \end{itemize}

    \item \textbf{Tests based on limits:}
        \begin{itemize}
            \item $n$-th term test
            \begin{itemize}
                \item \textit{Condition:} Can always be used

                \item \textit{Conclusion:} If $\displaystyle\lim_{k\to \infty} a_k \neq 0$ the series diverges
            \end{itemize}

            \item Finding closed form for partial sums $s_n$
            \begin{itemize}
                \item \textit{Condition:} In principle can always be used (in practice mostly helps in theoretical settings)

                \item \textit{Conclusion:} If $\displaystyle \lim_{k \to \infty} s_n$ exists the series converges; if the limit does not exist the series diverges
            \end{itemize}

            \item Ratio test
            \begin{itemize}
                \item \textit{Condition:} Can always be applied as long as $a_k \neq 0$

                \item \textit{Conclusion:} If $L = \displaystyle\lim_{k \to \infty} |a_{k+1}/a_k| < 1$ the series converges absolutely. If $L>1$ the series diverges
            \end{itemize}

            
        \end{itemize}
    

    \item \textbf{Comparison tests:}
        \begin{itemize}
            \item Direct comparison test
            \begin{itemize}
                \item \textit{Condition:} requires the terms of the series to be non-negative

                \item \textit{Conclusion:} Assume $0 \leq a_k \leq b_k$. If $\sum a_k$ diverges, then $\sum b_k$ diverges. If $\sum b_k$ converges, then $\sum a_k$ converges.
            \end{itemize}

            \item Integral comparison test
            \begin{itemize}
                \item \textit{Condition:} requires $a_k = f(k)$ where $f(x)$ is a non-negative, decreasing function defined on $[a,\infty)$

                \item \textit{Conclusion:} If $\displaystyle\int_a^\infty f(x) ~dx$ converges/diverges, then $\displaystyle\sum_{k=a}^\infty a_k$ converges/diverges as well
            \end{itemize}

            \item Asymptotic comparison test
            \begin{itemize}
                \item \textit{Condition:} requires terms to be non-negative

                \item \textit{Conclusion:} Assume $0 \leq a_k \sim b_k$. If $\sum b_k$ converges/diverges, then $\sum a_k$ covnerges diverges as well
            \end{itemize}

            \item Absolute convergence
            \begin{itemize}
                \item \textit{Condition:} can always be used

                \item \textit{Conclusion:} if $\sum |a_k|$ converges, then $\sum a_k$ converges absolutely (and, in particular, converges)
            \end{itemize}

            
        \end{itemize}
\end{enumerate}

        
    \end{solution}
        

    \item 
    \begin{problem}[2.5in]
    For each test above, find one series for which that test would be appropriate. Try to make a choice such that no other test would be as effective.    
    \end{problem}
    

    \item
    \begin{problem}[1.5in]
    Come up with a strategy for deciding what order of tests you should try when you encounter a new series.
    \end{problem}
\end{enumerate}



\pspace{\newpage}
\begin{problemonly}
\begin{framed}
\textbf{Theorem:} A function $f(x)$ is \textit{equal} to its Taylor series on its interval of convergence if and only if the exact error from its $n$th degree Taylor approximation approaches 0 as $n$ approaches $\infty$. That is:
$$f(x)=\text{ Taylor series of } f(x) \qquad \Longleftrightarrow \qquad \lim_{n\to \infty} \left|f(x)-P_n(x)\right|=0$$
for all $x$ in the interval of convergence.
\end{framed}
\end{problemonly}


\item 
  \note{I want to make a big deal of this exercise and go through it carefully to
  make the point that the Remainder Theorem helps us resolve the big questions that have bugged us for weeks (e.g., how good of an approximation is $P_{n}(x)$, and does a Taylor series converge to the function it was intended to represent).}

  Let $f(x) = e^{x}$.  
  \begin{enumerate}
     \item 
     \begin{problem}[1in]
       Write down the Taylor series for $f(x)$ centered at $0$.  What is its interval of convergence?
     \end{problem}

     \begin{solution}
       In Problem Set 13 \#1, you computed this Taylor series.  Because $f^{(k)}(x) = e^{x}$ for each $k \geq 0$,
       we have $f^{(k)}(0) = 1$.  It follows that the Taylor series for $f$ centered at $0$ is given by
       $$
	 \sum_{k = 0}^{\infty} \frac{1}{k!} x^{k}.
       $$
       By the Ratio Test, this series converges for all $x$.
     \end{solution}


     \item 
     \begin{problem}[1in]
       Give an upper bound on how much error is committed in using
       $P_{10}(3)$ to approximate $f(3) = e^{3}$.
     \end{problem}

     \begin{solution}
	Using the notation in the statement of the Remainder Theorem, we'll apply the theorem using
        $n = 10$, $d = 3$, and $c = 0$.  To do so, we need to supply $M_{(11)}$, an upper bound on
        $|f^{(11)}(x)|$ for $x$ the interval $[-3,3]$.  Since $f^{(11)}(x) = e^{x}$ is positive and increasing,
        we see that $|f^{(11)}(x)|$ achieves a maximum at the right endpoint of the domain $[-3,3]$.
        Therefore, a safe choice for $M_{(11)}$ would be
        $$
	  M_{(11)} = e^{3}.
        $$
        The Remainder Theorem tells us that
        $$
          |f(3) - P_{10}(3)| \leq \frac{e^{3}}{11!}\cdot (3-0)^{11} \leq 0.08914.
        $$
        In fact, it also tells us that $|f(x) - P_{10}(x)| \leq 0.08914$ {\em for all $x$}
        in the domain $[-3, 3]$.
     \end{solution}

     \item 
     \begin{problem}[1in]
       How large should $n$ be in order to guarantee that $|f(3) - P_{n}(3)| \leq 10^{-10}$?    
     \end{problem}

     \begin{solution}
        The Remainder Theorem tells us that
        $$
          |f(3) - P_{n}(3)| \leq \frac{M_{(n+1)}}{(n+1)!}\cdot |3-0|^{n+1},
        $$
        where $M_{(n+1)}$ is an upper bound on $|f^{(n+1)}(x)|$ for $x$ in the
        interval $[-3,3]$.  Since $f^{(n+1)} = e^{x}$, the same reasoning we
        used in the preceding part ensures that a safe choice for $M_{(n+1)}$ is $e^{3}$.
        Therefore,
        $$
          |f(3) - P_{n}(3)| \leq \frac{e^{3}}{(n+1)!}\cdot 3^{n+1}.
        $$
        If we can produce an $n$ for which
        $$
          \frac{e^{3}}{(n+1)!}\cdot 3^{n+1} \leq 10^{-10},
        $$
        then we'll know that $|f(3) - P_{n}(3)| \leq 10^{-10}$ as well.  You're welcome
        to try finding such an $n$ by pen-and-paper ``practical" estimates.  If you use 
        the WolframAlpha website to experiment with various choices of $n$, you'll find
        that $n = 22$ is the smallest choice that works.  
     \end{solution}

     \item 
     \begin{problem}[1in]
       Use the Remainder Theorem to show that the series you wrote down in Part (a)
       converges to $e^{x}$ for all $x$ in $[-3,3]$.  
     \end{problem}

     \begin{solution}
        In Part (c), we noted that
        $$
          |f(x) - P_{n}(x)| \leq \frac{e^{3} \cdot 3^{n+1}}{(n+1)!}
        $$
        for every $x$ in the interval $[-3,3]$.  Note that the series in Part (a)
        is a limit of Taylor polynomials, namely $\displaystyle \lim_{n \to \infty} P_{n}(x)$.
        If we were to let $n \to \infty$,
        the right-hand side of the above inequality will approach zero because
        $3^{n+1} \ll (n+1)!$.  It follows that the Taylor series is an {\em exact}
        representation of $e^{x}$ throughout the interval $[-3,3]$.
     \end{solution}

     \item 
     \begin{problem}[1in]
       Generalize Part (d) to show that the series you wrote down in Part (a) converges to $e^{x}$
       for all real $x$.  That is, the series really {\em is} a power series representation for $e^{x}$.
     \end{problem}

     \begin{solution}
        There was nothing special about the choice $d = 3$ in the above calculations.
        Suppose $d > 0$ and retrace the steps in the preceding parts of this problem. 
        You'll find that $M_{(n+1)} = e^{d}$ is an upper bound on $|f^{(n+1)}(x)|$
        for $x$ in the interval $[-d, d]$.  The Remainder Theorem ensures that
        $$
          |f(x) - P_{n}(x)| \leq \frac{e^{d} \cdot d^{n+1}}{(n+1)!}
        $$
        for every $x$ in the interval $[-d,d]$.  As before, note that $d^{n+1} \ll (n+1)!$,
        which means that $|f(x) - P_{n}(x)| \to 0$ for all $x$ in $[-d,d]$ as $n \to \infty$.
        Because $d$ was arbitrary, we conclude that the Taylor series converges to $e^x$
        for all real $x$.  That is, the series really {\em is} a power series representation for $e^x$.
     \end{solution}


  \end{enumerate}
  
  \pspace{\newpage}

 \item  \label{alt-series-bound}
 Consider the series
 $$
   \sum_{k = 10}^{\infty} (-1)^{k} \frac{1}{\ln(\ln k)}.
 $$

 \begin{enumerate}
   \item  
   \begin{problem}[1.5in]
     Explain why this series converges.
   \end{problem}

   \begin{solution}
     The Alternating Series Test applies.  To see why, note that
     $$
       |a_{k+1}| = \frac{1}{\ln(\ln(k+1))}  \; \leq \; \frac{1}{\ln(\ln k)} \; = \; |a_{k}|
     $$
     for each $k \geq 10$.  Moreover, $\displaystyle \lim_{k \to \infty} a_{k} = 0$.
     Therefore, the series converges by the Alternating Series Test.
     (As an exercise, you should explain why the series does NOT converge absolutely.)
   \end{solution}

   \item
   \begin{problem}[1.5in]
     Let $S$ denote the number that the above series converges to, and for each $n \geq 10$, let
     $$
       s_{n} = \sum_{k = 10}^{n} (-1)^{k} \frac{1}{\ln(\ln k)}.
     $$
     Give an example of an $n$ for which $s_{n}$ is definitely within $0.01$ of $S$.
   \end{problem}

   \begin{solution}
     Because this series satisfies the conditions of the Alternating Series Test, the
     associated error bound applies.  That is,
     $$
       |s_{n} - S| \; \leq \; \frac{1}{\ln(\ln(n+1))}.
     $$
     In order to guarantee that the partial sum $s_{n}$ is within $0.01$ of $S$, we
     can insist that
     $$
       \frac{1}{\ln(\ln(n+1))} \; \leq \; 0.01.
     $$
     Equivalently, we need $\ln(\ln(n+1)) \geq 100$, which means that
     $$
       n \geq e^{e^{100}}.
     $$
     So we can use any integer choice of $n$ that's larger than $\displaystyle e^{e^{100}}$, for example $3^{3^{100}}$.
     This is an absurdly large number (it has more than $10^{43}$ digits).  It would be
     impossible to compute such a partial sum, even if all of the computers on Earth had
     been working in tandem since the dawn of the universe.
   \end{solution}
   \end{enumerate} 
   
\pspace{\newpage}


    


\item 
For each series below, describe a strategy for finding an upper bound for the error associated to the sum of the first 50 terms, $s_{50}$:

\begin{enumerate}
    \item 
    \begin{problem}[\vfill]
        $\displaystyle \sum_{k=1}^\infty \frac{3^k}{9^k+1}$
    \end{problem}

    \begin{solution}
        The error associated to the approximation $s_{50}$ is the sum $\displaystyle \sum_{k=51}^\infty \frac{3^k}{9^k+1}$. We can upper bound that sum by the geometric series $\displaystyle \sum_{k=51}^\infty \frac{3^k}{9^k} = \sum_{k=51}^\infty \frac{1}{3^k}$ which converges to $\frac{1/3^{51}}{1- 1/3}$. Thus 
        \[\textrm{Error} \leq \frac{1}{2\cdot 3^{50}}\]
    \end{solution}

    \item 
    \begin{problem}[\vfill]
        $\displaystyle \sum_{k=1}^\infty \frac{1}{k^3}$
    \end{problem}

    \begin{solution}
        The error associated to the approximation $s_{50}$ is the sum $\displaystyle \sum_{k=51}^\infty \frac{1}{k^3}$. We can upper bound that sum by the integral $\displaystyle\int_{50}^{\infty} \frac{1}{x^3} ~dx$, which can be computed and results in $\frac{1}{2 \cdot 50^2}$. Thus,
        \[\textrm{Error} \leq \frac{1}{2 \cdot 50^2}\]
    \end{solution}
    
\end{enumerate}







\pspace{\newpage}

\item 
 Your roommate is working on some series practice problems.  In each part, decide
 whether their solution is reasoned correctly.  If it isn't, explain
 exactly what is wrong with their reasoning and provide a correct argument to conclude convergence or divergence of the given series.
 
 \begin{enumerate}
 \item
   \begin{problem}
     \textbf{Problem:} Decide whether $\displaystyle \sum_{k=10}^\infty 
     \frac{k^2 + 2^k}{k! - k}$ converges.

     \textbf{Your roommate's answer:} $\displaystyle \frac{k^2 + 2^k}{k! - k} \sim
     \frac{2^k}{k!}$ because $2^k \gg k^2$ and $k! \gg k$.  I can use the
     Asymptotic Comparison Test to say that $\displaystyle \sum \frac{k^2
       + 2^k}{k! - k}$ behaves the same way as $\displaystyle \sum
     \frac{2^k}{k!}$ since $\displaystyle \frac{2^k}{k!} > 0$.  I know
     that $2^k \ll k!$, so $\displaystyle \sum \frac{2^k}{k!}$ converges
     because the numerator $\ll$ the denominator.  Therefore,
     $\displaystyle \sum_{k=10}^\infty \frac{k^2 + 2^k}{k! - k}$
     converges as well.
   \end{problem}

   \begin{solution}
     Your roommate's reasoning is invalid.  They are correct that $\displaystyle
     \frac{k^2 + 2^k}{k! - k} \sim \frac{2^k}{k!}$ and that $2^k \ll k!$.
     However, the reasoning that $\sum \frac{2^k}{k!}$ converges because
     numerator $\ll$ denominator is incorrect.

     Here's an example which demonstrates that this reasoning is invalid:
     $1 \ll k$, but $\sum \frac{1}{k}$ still diverges.

    We could use the asymptotic comparison test, but apply the ratio test to the simplified series $\displaystyle \sum_{k=10}^\infty 
     \frac{2^k}{k!}$. Notice that

     $$\lim_{k\to \infty} \left|\frac{a_{k+1}}{a_k}\right| = \lim_{k\to \infty} \left|\frac{2^{k+1}}{(k+1)!}\frac{k!}{2^k}\right| = \lim_{k\to \infty} \frac{2}{k+1} = 0 <1$$

     So the series $\displaystyle \sum_{k=10}^\infty \frac{2^k}{k!}$ converges by the ratio test and thus $\displaystyle \sum_{k=10}^\infty \frac{k^2 + 2^k}{k! - k}$ converges by the asymptotic comparison test.
     
   \end{solution}




 \item
   \begin{problem}
     \textbf{Problem:} Decide whether $\displaystyle \sum_{k=10}^\infty
     \frac{1}{k \cdot 4^k}$ converges.

     \textbf{Your roommate's answer:} Since $k \ll 4^k$, $\displaystyle \frac{1}{k
       \cdot 4^k} \sim \frac{1}{4^k}$.  We know that $\sum \frac{1}{4^k}$
     converges because it's a geometric series with common ratio $r=
     \frac{1}{4}$, and $\left|\frac{1}{4} \right| < 1$.  Since $\sum
     \frac{1}{4^k}$ has positive terms, we can apply the Asymptotic
     Comparison Test to conclude that $\sum \frac{1}{k \cdot 4^k}$ also
     converges. 
   \end{problem}

   \begin{solution}
     Your roommate is correct that $k \ll 4^k$, but it is \textbf{not} true that
     $\displaystyle \frac{1}{k \cdot 4^k} \sim \frac{1}{4^k}$.  (This is
     a common mistake: a \textit{sum} is asymptotic to its largest term,
     but the same is not true for products.)

     A valid argument would be to use the direct comparison test. Notice that 

     $$0 \leq \frac{1}{k 4^k} \leq \frac{1}{4^k}$$

     We know that $\sum_{k=10}^\infty \frac{1}{4^k}$ is a geometric series with $|r| = 1/4 <1$, so it converges. Therefore, by the direct comparison test the original series converges as well.
   \end{solution}

 \item
   \begin{problem}
     \textbf{Problem:} Decide whether $\displaystyle \sum_{k=2}^\infty
     \frac{1}{k (\ln k)^2}$ converges.

     \textbf{Your roommate's answer:} $(\ln k)^2 \gg 1$, so $\displaystyle
     \frac{1}{k (\ln k)^2} \ll \frac{1}{k}$.  This means that
     $\displaystyle \frac{1}{k (\ln k)^2}$ goes to 0 much more quickly
     than $\displaystyle \frac{1}{k}$ does.  Since $\displaystyle \sum
     \frac{1}{k}$ just barely diverges, $\displaystyle \sum_{k=2}^\infty
     \frac{1}{k (\ln k)^2}$ must converge.
   \end{problem}

   \begin{solution}
     Your roommate is correct that $\displaystyle \frac{1}{k (\ln k)^2} \ll
     \frac{1}{k}$.  However, they have fallen into the ``harmonic series as
     dividing line'' trap (see Problem Set 16, {\#2}) of thinking
     that things which are smaller than the harmonic series must
     converge.  Remember: If you're trying to show that a series
     converges by comparing it with another series, you must compare it
     with a convergent series.

     We can apply the integral comparison test. Let us look at the improper integral
     \begin{align*}
     \int_2^\infty \frac{1}{x (\ln x)^2}~dx & = \lim_{b \to \infty} \int_2^b \frac{1}{x (\ln x)^2}~dx =  \lim_{b \to \infty} \int_{x=2}^{x=b} \frac{1}{u^2}~du\\
     & = \lim_{b \to \infty} -\left.\frac{1}{(\ln x)^2} \right|_{x=2}^{x=b} = \lim_{b \to \infty} \left(-\frac{1}{(\ln b)^2}+\frac{1}{(\ln 2)^2}\right) = \frac{1}{(\ln 2)^2}
     \end{align*}

     Since the integral above converges, by the integral comparsion test we know that the series $\sum_{k=2}^\infty
     \frac{1}{k (\ln k)^2}$ converges as well.
    \end{solution}


 \item
   \begin{problem}
     \textbf{Problem:} Decide whether $\displaystyle \sum_{k=1743}^\infty
     \frac{0.95^k + k^{1.05}}{k \sqrt{k}}$ converges.

     \textbf{Your roommate's answer:} We can rewrite $k \sqrt{k}$ as $k^{1.5}$.
     Since $0.95^k \ll k^{1.05}$, $\displaystyle \frac{0.95^k +
       k^{1.05}}{k^{1.5}} \sim \frac{k^{1.05}}{k^{1.5}} =
     \frac{k^{1.05}}{k^{1.5}} = \frac{1}{k^{0.45}}$.  We know that
     $\displaystyle \sum \frac{1}{k^{0.45}}$ diverges because it's a
     $p$-series with $p = 0.45 \leq 1$.  Since $\displaystyle \sum
     \frac{1}{k^{0.45}}$ has positive terms, we can apply the Asymptotic
     Comparison Test to conclude that $\displaystyle \sum \frac{0.95^k +
       k^{1.05}}{k \sqrt{k}}$ also diverges.
   \end{problem}

   \begin{solution}
     Your roommate is completely correct.
   \end{solution}
   
 \end{enumerate}

 \begin{problemonly}
   \textit{Hint:} Only one of your roommate's answers is completely correct.
 \end{problemonly}


 \end{enumerate}



\end{document}